% Figaro Paper J: Air Service as Coordinated Resource Markets
% Industrial / Systems Engineering register.
% Audience: industrial-engineering and systems-engineering scholars,
% supply-chain-coordination researchers, scheduled-service operations
% practitioners.
% Preprint, May 2026

\documentclass[11pt,a4paper]{article}

\usepackage[utf8]{inputenc}
\usepackage[T1]{fontenc}
\usepackage{amsmath,amssymb,amsthm}
\usepackage{mathtools}
\usepackage{booktabs}
\usepackage{graphicx}
\usepackage{hyperref}
\usepackage{cleveref}
\usepackage{enumitem}
\usepackage{xcolor}
\usepackage[margin=1in]{geometry}
\usepackage{natbib}
\bibliographystyle{plainnat}

\newtheorem{theorem}{Theorem}[section]
\newtheorem{proposition}[theorem]{Proposition}
\newtheorem{corollary}[theorem]{Corollary}
\theoremstyle{definition}
\newtheorem{definition}[theorem]{Definition}
\newtheorem{example}[theorem]{Example}
\theoremstyle{remark}
\newtheorem{remark}[theorem]{Remark}

\newcommand{\figaro}{\textsc{Figaro}}
\newcommand{\pay}{P}
\newcommand{\cumval}{G}

\title{%
  Air Service as Coordinated Resource Markets:\\
  An Industrial-Engineering Reading of the Bonded Architecture%
}

\author{
  Alessandro Daliana
}

\date{May 2026}

\begin{document}

\maketitle

% ════════════════════════════════════════════════════════════════════
\begin{abstract}
Air service is coordination across a set of resource markets:
crew-hours and certifications, aircraft-time, fuel volume, gate
slots, catering counts, maintenance hours, ground handling.
Each resource provider is an off-chain real-world asset (RWA)
--- a pilot, an airframe, a fuel inventory, a security
checkpoint --- whose on-chain participation is mediated by a
wallet. The wallet holds the asset's earnings as token balances
and points to the asset's off-chain credentials through NFTs (a
flight-crew certification, an airworthiness certificate, an
operating permit). A human or autonomous agent operates the
wallet on the asset's behalf. The wallet is an active counterparty
in the assembly iff the asset is alive and its credentials are
current. Public-authority services participate on the same
footing as commercial providers. Bonded commitments
in the Figaro kernel coordinate these wallets directly: the
buyer wallet (the passenger) posts $2P_i$ at each commit, and
each provider wallet posts $2G$ against the cumulative process
value at the moment it commits.

Two architectural properties carry the contribution.
\emph{Asymmetric bonding within a single process}: the passenger
is the \texttt{rootBuyer} of every order in the process; each
provider wallet bonds against the cumulative-value accumulator
$G$ at its own commit; the kernel enforces a single
\texttt{rootBuyer} across every order. \emph{Buyer-initiated
atomic resolution}: the passenger settles the entire process at
destination as a single action; every active order in the process
resolves together or none does.

Two consequences follow. First, the resource markets are
directly accessible to the passenger through bonded commerce.
Each provider wallet's interests align with the passenger's at
the bonded commit, mediated by no third party. Second, when
operational disruption occurs the seller cohort faces direct
economic incentive --- under locked bonds across every edge --- to
determine and pay compensation directly to the buyer, before any
external mechanism (specialized contracts, dispute forums,
insurance, regulators) engages. Cascading-delay reduction
follows as a consequence of bond posture and atomic resolution;
it is a property the architecture produces, not the property the
paper is for.
\end{abstract}

\medskip
\noindent\textbf{Keywords:}
scheduled-service coordination, supply-chain coordination,
process modeling, resource markets, weakest-link coordination,
industrial engineering

% ════════════════════════════════════════════════════════════════════
\section{Introduction}
\label{sec:intro}

Air service is coordination across a set of resource markets.
The journey a passenger purchases is the joint output of a
flight crew (whose certifications and duty-time bands are
verified at agreement signing), an aircraft (whose airworthiness
certificate is verified the same way), a quantity of jet fuel,
slot allocations at origin and destination gates, per-passenger
catering, a share of the day's maintenance allocation for the
operating airframe, and ground handling for boarding and
baggage. Each is a distinct resource the journey draws on; each
has natural units of measure already in industry use --- crew-hours,
type ratings, fuel volume in gallons or kilograms, gate-slot time
bands, meal counts, maintenance hours, baggage counts. Public
services are further inputs the flight consumes: air-traffic
control, security screening, the federal aviation system itself.
These public-authority services are real-world assets with
commercial value in the same sense the catering supply or the
fuel inventory is, and their wallets participate in the bonded
process on the same footing.

Each provider in the assembly is an RWA-as-wallet. The
\emph{asset} is the off-chain real-world asset whose
participation produces value --- the airframe that flies, the
labor capacity of the pilot, the screening service the security
authority provides. The \emph{wallet} is the on-chain
representation of that participation: an address holding the
asset's earnings as token balances, plus NFTs that point to the
off-chain credentials the asset's role requires (the pilot's
flight-crew certification, the aircraft's airworthiness
certificate, the maintenance provider's authority-issued
authorization). The \emph{operator} is the human or autonomous
agent who controls the wallet's signing key on the asset's
behalf. The wallet is an active counterparty in the assembly iff
the asset is alive and its credentials are current; that
conditional is enforced operationally by the operator and at
agreement signing by whichever credential schemas the assembly
designer composes.

Each provider wallet sustains its participation the way a node
on a blockchain network sustains its participation. A node earns
fees and rewards from the network and uses them to cover its
off-chain operating expenses (machines, electricity, HVAC, rent);
when the receipts no longer cover the costs, the operator stops
running the node and the node drops out of the network. An
RWA-as-wallet operates on the same logic in its market: receipts
from bonded processes must cover the asset's off-chain operating
expenses (fuel, parts, labor, premises, regulatory fees, capital
amortization), or the operator stops bonding the wallet into
processes and the wallet drops out of the market it serves. This
holds for the catering wallet and for the security-screening
wallet alike: the kernel does not distinguish between a
"market-traded fee" and a "regulated fee." Both are payments to
service-providing wallets that need receipts sufficient to
sustain operation. The schemas the assembly uses to verify
credentials at agreement signing are generic across vehicle types
--- the same credential and schedule schemas serve aircraft,
ships, trucks, and other vehicles, with per-vehicle personalization
in the schema content and validator-contract configuration.

The bonded settlement primitive coordinates across these wallets
without a coordinating party. Enforcement of bilateral commitments
is self-enforcing through asymmetric bonding; multi-party
coordination is resolvable from a single signature through atomic
resolution. The journey is composed as a bonded process across
those resource markets; each provider wallet bonds against the
passenger's process; the passenger settles the entire process at
destination through one atomic \texttt{resolveProcess} call.

\paragraph{Paper organization.}
\Cref{sec:primitive} summarizes the settlement primitive.
\Cref{sec:assembly} develops the architectural properties applied
to air-service coordination: per-edge asymmetric bonding,
buyer-initiated atomic resolution, the two consequences for the
assembly, the kernel / composability-stack distinction, and where
specialized mechanism contracts can be developed for specific
resource markets. \Cref{sec:disruption} treats disruption
resolution: the seller-cohort compensation mechanism that arises
from buyer-dominance under bond pressure (the primary mechanism),
illustrated through a worked bond-posture example and through the
crew-allocation case, with the fallback mechanisms (specialized
contracts, coordination tokens, off-chain dispute forums,
insurance, regulators) treated downstream. \Cref{sec:schemas}
treats the schema requirements: a sister-schema pair for
schedule binding, with the rest of the assembly satisfied by
existing infrastructure that is generic across vehicle types.
\Cref{sec:comparison} treats what the architecture preserves
(safety regulation, crew-labor protections, insurance,
competition law) and what it changes (the cost-flow path).
\Cref{sec:scope} states the scope conditions.
\Cref{sec:conclusion} concludes.

% ════════════════════════════════════════════════════════════════════
\section{The Settlement Primitive}
\label{sec:primitive}

The bonded primitive in \texttt{FigaroCore} has two mechanisms.
We summarize the features the present argument uses.

\paragraph{Asymmetric bonding (Mechanism~1).}
For a transaction with payment $\pay > 0$ and cumulative upstream
value $\cumval \geq \pay$, the buyer locks $2\pay$ and the seller
locks $2\cumval$. Cooperation is the unique profile surviving
iterated elimination of weakly dominated strategies. The
mechanism scales to $N$-party processes through progressive
collateralization: each seller bonds against the cumulative value
of all upstream commitments, producing a mesh of independently
secured bilateral edges, each of which carries its own
equilibrium.

\paragraph{Buyer dominance with atomic resolution (Mechanism~2).}
Only the root buyer can trigger resolution, and resolution
settles all active orders in the process simultaneously or not at
all. The all-or-nothing rule induces a weakest-link subgame among
sellers: each seller's payout depends on universal cooperation
across the process. The architecture produces cooperation
pressure of magnitude $\pay_i + 2\cumval_i$ on every other seller
from each seller's perspective, without explicit communication,
governance, or managerial layer.

\paragraph{What this implies for air-service coordination.}
The bonded primitive's two mechanisms supply two coordinative
properties for air service. Mechanism~1 supplies a bilateral
commitment whose settlement does not require either party's home
jurisdiction's contract-law regime to enforce: settlement is the
on-chain release of locked bonds, mediated by no court. The
off-chain agreement remains binding under whatever law the
parties select for the substantive commitment, but the kernel's
settlement guarantee is jurisdiction-agnostic. Mechanism~2
supplies the atomic resolution that propagates a single
commit-or-don't decision through the entire process; this is
what makes operational damages flow proportionally to the unit
that caused them rather than being absorbed at whichever contract
has the loosest cap. Neither mechanism is novel to the
air-service domain; both are general-purpose settlement
properties whose application here is the present paper's
contribution.

% ════════════════════════════════════════════════════════════════════
\section{The Bonded Architecture Applied to Air Service}
\label{sec:assembly}

The substantive contribution of the bonded primitive to
air-service coordination operates at the edge of the bonded
process, not at any particular position within it. Two
architectural properties supply the contribution.

\paragraph{Asymmetric bonding within a single process.}
The kernel enforces a strong structural rule: every order in a
process has the same buyer (the \texttt{rootBuyer}). The
passenger wallet commits directly to each provider wallet. The
first commit creates the process under the passenger's address;
every subsequent commit in the same process extends the
cumulative-value accumulator $G$ by the new payment. At each
commit the passenger posts $2P_i$ for that specific payment, and
the seller posts $2G_{\text{at commit time}}$ --- the full
cumulative process value at the moment they commit. Bond posture
is asymmetric: the passenger's per-commit bond scales with the
immediate payment, while each seller's bond scales with the
cumulative value of every upstream commit in the process.
Sellers committing later post correspondingly larger bonds; the
last seller bonds against the entire ticket value.

\paragraph{Buyer-initiated atomic resolution.}
The passenger calls \texttt{resolveProcess}, and the entire
process settles simultaneously. Every active order in the process
is resolved together or none is. The passenger's signature is the
only signature required for settlement, and atomicity guarantees
that no subset of orders settles without the rest.

\paragraph{Two consequences for the air-service assembly.}
The two properties together produce two consequences.

First, \emph{the resource markets named in \Cref{sec:intro}
become directly accessible to the passenger through bonded
commerce.} Each provider wallet --- the crew member wallet, the
aircraft wallet, the fuel-supplier wallet, the gate-operator
wallet, the caterer wallet, the maintenance-provider wallet, the
ground-handling wallet --- is a seller bonded against the
passenger's process, with payout depending on the passenger's
atomic resolution. Each provider's interests align with the
passenger's at the bonded commit, mediated by no third party.
We treat one alignment case in detail in \Cref{sec:disruption}:
the crew-allocation arrangement.

Second, \emph{damage flow tracks bond posture across the
assembly's commits rather than the loosest contractual liability
cap.} Each seller carries asymmetric bond exposure proportional
to the cumulative process value at their commit; the seller-side
total grows super-linearly with the number of commits while the
passenger's posture stays at $2 P_{\text{ticket}}$. The
atomic-resolution rule operates on this set of bonds. When
something goes wrong, every seller's bond is locked across every
commit --- which is the structural pressure behind the
disruption-resolution mechanism we develop in
\Cref{sec:disruption}.

\paragraph{What the kernel constrains; what the composability stack carries.}
The kernel reads only per-order facts: buyer (which must be the
\texttt{rootBuyer} in every order in the process), seller,
payment, cumulative-value snapshot at commit, and the
atomic-resolution rule covering the process. The kernel-level
shape of a single process is therefore tightly constrained: one
buyer, one or more sellers, all under one cumulative accumulator,
all resolving atomically. Deeper compositions across multiple
processes (a seller in process $A$ becoming the
\texttt{rootBuyer} of process $B$ to coordinate their own
sub-suppliers) are admissible at the kernel layer but resolve
process-by-process, not across.

Above the kernel, the protocol-and-runtime composability stack
attaches schemas to specific orders (the
\texttt{scheduled-departure-v1} clause binds to whichever order
the parties want to subject to schedule-tolerance enforcement;
\texttt{handoff-v1} attestations attach to operational-handoff
orders), mechanism contracts operate on specific order patterns,
assembly specifications declare which schemas attach where and
which view definitions render which surfaces, and UIs render the
process per the assembly's declarations. An assembly's structure,
once specified, is a fixed verifiable artifact bound to its full
composability stack. The bonded-architecture argument of this
paper applies to whichever resource markets the designer composes
under one process, with the kernel-enforced structural rule
holding throughout.

\paragraph{Specialized mechanism contracts per resource market.}
Each resource market admits its own market mechanism. The
existing \texttt{DutchAuction} contract supplies one such
mechanism for cases where descending-price coordination is
appropriate (fuel procurement is a candidate where the supplier
market admits competition). Bilateral commerce-v1 commits cover
the cases where the parties contract directly. Additional
mechanism contracts can be developed under the protocol-extension
discipline for resources whose market mechanics deserve their own
primitive: a slot-allocation mechanism for gate-time bidding, a
crew-dispatch mechanism that respects type-rating and duty-time
constraints, a fuel-pool mechanism for shared ramp-side
procurement. The mechanism-contract space sits at the runtime
tier; the kernel constraint that every order in a process shares
one \texttt{rootBuyer} (the passenger) holds across whichever
mechanism contracts the assembly composes.

% ════════════════════════════════════════════════════════════════════
\section{Disruption Resolution Under Buyer Dominance}
\label{sec:disruption}

We turn to the operational question: when something goes wrong
--- a delay, a slip, a service failure --- how is the disruption
resolved under the bonded architecture? The substantive answer is
that Mechanism~2 of the kernel (buyer dominance with atomic
resolution) creates structural pressure for the seller cohort to
determine and pay compensation themselves, directly to the buyer,
before any external mechanism engages. This is the primary
damage-management mechanism in the architecture. We treat it
first, place a bond-posture worked example, develop the
crew-allocation case as a concrete illustration, note the
IROps-scale extension, and then position the fallback mechanisms
in their proper register downstream.

\subsection{The primary mechanism: seller-cohort compensation under bond pressure}

\paragraph{Bond posture per commit.}
Every order in a process shares the same \texttt{rootBuyer}; the
passenger as that \texttt{rootBuyer} posts $2P_i$ at each commit
to each resource provider, and each resource provider posts
$2G_{\text{at commit time}}$ where $G$ is the cumulative
process value at the moment of that commit. Sellers committing
later in the sequence post correspondingly larger bonds.
Aggregated across an $N$-seller process, the bonded exposure of
the seller cohort scales super-linearly with the number of
resource markets the assembly draws on.

\paragraph{A bond-posture worked example.}
Consider a stylized domestic short-haul passenger journey
totalling $\$250$. The kernel enforces $\texttt{buyer} =
\texttt{rootBuyer}$ in every order in a process, so the passenger
commits directly to each provider wallet. The ten commits below
are all part of one process under one $\texttt{rootBuyer}$ (the
passenger). The kernel maintains a single process-wide
cumulative-value accumulator $G$ that grows monotonically as each
commit lands ($G \leftarrow G + P_i$). At every commit the buyer
posts $2P_i$ and the seller posts $2G_{\text{at commit time}}$.
The asymmetry shows up across the sequence: the passenger's
buyer-side bond at each commit scales only with the immediate
payment to that seller, while each seller's bond scales with the
full cumulative value of every upstream commit in the process.
Sellers committing later post correspondingly larger bonds; the
last seller bonds against the entire ticket value. This is
progressive collateralization in operation.
\Cref{tab:bond-posture} traces the sequence.

\begin{table}[h]
\centering
\small
\begin{tabular}{lrrrr}
\toprule
Seller (commit sequence) & $P_i$ & $G$ after commit & Passenger bond $2P_i$ & Seller bond $2G$ \\
\midrule
Aircraft wallet (root) & $\$65$ & $\$65$ & $\$130$ & $\$130$ \\
Fuel-supplier wallet & $\$45$ & $\$110$ & $\$90$ & $\$220$ \\
Crew member wallet & $\$40$ & $\$150$ & $\$80$ & $\$300$ \\
Maintenance-provider wallet & $\$15$ & $\$165$ & $\$30$ & $\$330$ \\
Catering wallet & $\$5$ & $\$170$ & $\$10$ & $\$340$ \\
Ground-handling wallet & $\$15$ & $\$185$ & $\$30$ & $\$370$ \\
Origin airport authority wallet & $\$20$ & $\$205$ & $\$40$ & $\$410$ \\
Destination airport authority wallet & $\$20$ & $\$225$ & $\$40$ & $\$450$ \\
TSA wallet (security screening) & $\$6$ & $\$231$ & $\$12$ & $\$462$ \\
Federal aviation system wallet & $\$19$ & $\$250$ & $\$38$ & $\$500$ \\
\midrule
Total bond locked across the process & --- & --- & $\$500$ & $\$3{,}512$ \\
\bottomrule
\end{tabular}
\caption{Stylized per-commit bond posture for a domestic
short-haul passenger journey under one process. The passenger
commits directly to each provider wallet (the kernel enforces
buyer $= \texttt{rootBuyer}$ in every order). $G$ is the
process-wide cumulative-value accumulator, monotonically growing
with each commit. The passenger's bond at each commit is
proportional to the immediate payment; each seller's bond is
proportional to the cumulative process value at the moment that
seller commits. The last seller to commit therefore posts the
largest seller bond, $2G_{\text{final}} = 2 \times \$250 = \$500$.
Per-seller payments are illustrative; absolute magnitudes vary
with route, aircraft type, and operating context.}
\label{tab:bond-posture}
\end{table}

Two features of the table carry the structural argument.

\emph{The asymmetry between buyer and seller bonds.} The
passenger's aggregate locked bond is $\$500$, which equals
$2 \times P_{\text{ticket}}$ --- the standard buyer posture. The
seller cohort's aggregate locked bond is $\$3{,}512$, more than
seven times the passenger's posture, because each seller bonds
against the cumulative $G$ that has grown across the upstream
commits. This is not an accident of the example; it is the
asymmetric-bonding architecture in operation. As the assembly
adds more resource markets, the seller-side total grows
super-linearly with the number of commits while the passenger's
side stays at $2 P_{\text{ticket}}$.

\emph{The magnitude of the seller cohort's locked bond.} The
seller cohort's aggregate locked bond ($\$3{,}512$) is more than
fourteen times the ticket payment. A non-resolving passenger
holds that full pool locked across every commit in the process.
The pressure on the cohort to compensate the passenger and secure
resolution is structural in proportion that exceeds the ticket by
more than an order of magnitude.

The order in which sub-orders commit reorders which specific
seller posts the largest bond, but does not change the aggregate
seller cohort total: it depends only on the set of payments, not
the sequence. Whichever seller commits last bonds against the
full $G_{\text{final}} = P_{\text{ticket}}$.

\paragraph{Withholding resolution as the buyer's leverage.}
The passenger holds unilateral authority to call
\texttt{resolveProcess}. As long as the passenger does not call,
every seller's bond in the process remains locked. Each seller
faces forfeit of their full $2G_i$ on continued non-resolution.
The buyer's atomic-resolution authority is therefore not a remedy
in the conventional sense (it does not pay the buyer anything
beyond bond recovery); it is leverage. The buyer can hold the
entire bond pool locked by simply not signing.

\paragraph{Sellers determine and pay compensation themselves.}
When operational disruption occurs and the buyer is left worse
off than the agreement contemplated, the seller cohort faces
direct economic incentive to make the buyer whole enough that the
buyer chooses to resolve. The mechanism unfolds operationally as
follows.

Sellers approach the buyer with compensation offers. Compensation
takes whatever form the buyer accepts: cash transfer, future
service credit, accommodation or meal vouchers, refund of the
ticket payment, alternative routing, partial or full refund
combined with services. The currency of compensation is at the
parties' discretion --- the kernel does not constrain it.

Sellers negotiate among themselves about which seller bears what
fraction of the compensation. The unit whose performance
materially failed bears more. Operational parties around the
failure who could not have prevented it (the catering provider
who had no role in a fuel slip; the gate operator whose service
was unaffected) contribute less or not at all. The negotiation is
real but routine: each seller knows their own bond exposure,
their counterfactual payout under cooperative resolution, and
their willingness-to-pay against losing the bond entirely.

The buyer accepts the compensation that satisfies them and calls
\texttt{resolveProcess}. Bonds release across the process; every
seller recovers $2G_i + P_i$ less their share of the
compensation; the buyer recovers their bond plus the compensation
already paid. The settlement is bilateral and direct.

This happens fast. The negotiation occurs at the operational
boundary (the airport gate, the destination airport's customer
service desk, in some cases through the runtime UI's
disruption-coordination surface), not in a court filing six
months later. It happens before insurance claims, before
dispute-resolution forums, before regulatory engagement. Most
disruption resolves through this path because that is where the
bond pressure points.

\paragraph{The crew-allocation case.}
The mechanism is concrete in the crew market. The crew-allocation
operates as a direct market between the passenger and the crew
member. The crew member's wallet --- holding their certification
NFTs and operating-permit credentials, verified through the
relevant credential schemas at agreement signing --- bonds
against the passenger's process at the crew-edge. The crew
member's payout depends on the passenger's atomic resolution at
destination; the crew member's interests align with the
passenger's directly through the bonded edge. The allocation
discipline is whichever mechanism contract the assembly designer
composes for the crew market: a bilateral commerce-v1 commit, a
descending-price \texttt{DutchAuction}, or a future
crew-dispatch mechanism contract that respects type-rating and
duty-time constraints. Collective-bargaining institutions over
wage rates, duty-time protections, and seniority among crew
remain; they are labor-market institutions the bonded primitive
does not touch.

Under disruption --- say a crew member cannot fly --- the
crew-edge bond is locked alongside every other edge in the
process. The crew member negotiates compensation directly with
the passenger (find a substitute crew member at speed,
accommodate the passenger on the next flight, refund the leg, or
whatever the passenger accepts). Resolution follows compensation.
There is no intermediary in this negotiation; the architecture
does not contemplate one.

\paragraph{IROps-scale coordination.}
The single-flight example understates the operational complexity
of large-scale irregular operations (IROps). A hub disruption
produces simultaneous bonded processes across hundreds of
flights, with thousands of bonded edges across the resource-market
mesh. The compensation negotiation at IROps scale is operationally
heavier than the per-flight case: each affected passenger holds a
separate bonded process with its own seller cohort, and many
provider wallets (crew member wallets, aircraft wallets,
ground-handling wallets) appear in many processes at once. The
bonded architecture does not eliminate this coordination cost;
it routes the cost in a structurally different direction. Each
affected passenger faces a busy but bilaterally coordinated
seller cohort negotiating their compensation directly.
Operational tooling for the provider wallet's operator --- who
manages bonds locked across many simultaneous processes and must
coordinate compensation responses at scale --- is itself an
assembly-design question.

\paragraph{The weakest-link character of the pressure.}
Under buyer dominance with atomic resolution, each seller bears
direct economic interest of magnitude $P_i + 2G_i$ in every
other seller's performance, because each seller's payout depends
on universal cooperation across the process: cooperation is the
unique equilibrium surviving iterated elimination of weakly
dominated strategies, induced by atomic resolution operating on
the bonded mesh. The analytical content is the formal
characterization of the pressure; the operational content is
sellers paying compensation to buyers directly to secure
resolution. The same structural shape produces the joint-liability
dynamic in Grameen-style microfinance, in which group members
chip in to cover a struggling member's installment before the
lender forecloses, because each member's continued credit access
is exposed to every other member's repayment. The bonded
architecture's seller-cohort-compensates-buyer dynamic is the
same structural mechanism scaled with chain depth and operating
without the social-substrate prerequisites Grameen requires.
Buyer dominance with atomic resolution is in this sense a
\emph{social mechanism}: it produces social-coordination behavior
(peer pressure, cohort negotiation, burden-sharing, coverage of a
struggling counterpart) endogenously through the bond architecture,
with each provider wallet's continued participation in its market
making the negotiation rationally compelled rather than socially
expected.

\paragraph{Why this is structural, not procedural.}
The mechanism is not a courtesy that sellers may or may not
extend. It is a direct consequence of the bond posture under
atomic resolution. A seller who refuses to participate in
compensation negotiation watches their full bond forfeit
alongside every other seller's. The compensation mechanism
emerges from the architecture because the architecture has made
non-cooperation strictly worse than reasonable compensation for
every seller in the process.

\subsection{Fallback mechanisms}

The primary mechanism resolves most disruption directly. Fallback
mechanisms engage when the primary does not, or when the
disruption involves a class of issue the primary cannot address.
We list them in the order of operational engagement.

\emph{Specialized mechanism contracts (protocol layer).} Where
the assembly composes a specialized mechanism contract --- a
delay-compensation contract that disburses pre-agreed amounts
under specified conditions, a slot-reallocation mechanism, a
fuel-pool mechanism --- the contract handles its specific
coordination automatically. The specialized contract complements
the primary mechanism by making specific compensation patterns
automatic rather than negotiated.

\emph{Coordination tokens (runtime layer).} Where the assembly
composes around a coordination token, the substrate's operational
state registers as token value. Token holders bear operational
quality through value impact rather than through direct
compensation. Token-value reflection is a diffuse signal
operating in parallel with the per-edge compensation mechanism,
not a substitute for it.

\emph{Off-chain dispute resolution forums.} The off-chain
agreement named in the bonded commitment binds a dispute forum
(state court, international arbitration center, ODR platform,
Schelling-point juror system). The forum engages when the
seller-cohort compensation does not resolve the dispute ---
typically because the dispute is substantive in a way bond
pressure cannot address (duress, mistake, frustration,
illegality, public-policy concerns). The bonded record is the
evidentiary input the forum operates on: kernel byproducts
(commit and resolution events) supply the indisputable settlement
record, and schema-validated attestations (\texttt{handoff-v1}
timestamps, \texttt{courier-process-v1} stage events) supply
the substantive performance evidence.

\emph{Insurance compositions.} A passenger or provider wallet
who wants insurance against risks the bond pool cannot absorb
(catastrophic events, multi-day weather disruption, force-majeure
events) composes a parallel bonded process with an insurer. The
insurer's bond underwrites the risk; insurance settlement is its
own bonded event. Insurance composes with the bonded primitive
without modifying the kernel.

\emph{Regulators.} Aviation safety regulators, consumer protection
authorities, antitrust enforcement, and labor protection regimes
operate alongside the bonded architecture and engage in the
matters their authority covers. The architecture does not
displace them.

The substantive contribution of the bonded architecture to damage
management is that the primary mechanism --- seller cohort
negotiating and paying compensation directly to the buyer, before
external mechanisms engage --- is structural, immediate, and
operates without any of the fallbacks being needed in most cases.
The conventional regime's cascading-delay pathology --- damages
absorbed by whichever contract has the loosest liability cap,
with limited unit-level operational signal
\citep{lan_clarke_barnhart_2006} --- reduces under the bonded
architecture as a consequence of bond posture and atomic
resolution.

% ════════════════════════════════════════════════════════════════════
\section{Schemas}
\label{sec:schemas}

The assembly draws almost entirely on existing schema
infrastructure. A sister-schema pair is required for schedule
binding: \texttt{scheduled-departure-v1} and
\texttt{scheduled-arrival-v1}, so the resolution path can verify
whether the sellers delivered against the agreed departure and
arrival schedule respectively.

The schedule and credential schemas the assembly relies on are
generic across vehicle types: the same primitives serve aircraft,
ships, trucks, and other scheduled vehicles, with per-vehicle
personalization in the schema content (designator format,
authority issuing the credential, tolerance windows) and in
validator-contract configuration. The air-service application
below is one personalization of those primitives, not a
domain-specific architecture.

\paragraph{Existing schemas.}
\texttt{figaro-commerce-v1} carries the currency, payment, and
itemized line-items of every commitment in the assembly --- the
passenger ticket and every sub-procurement commitment.
\texttt{figaro-courier-process-v1} carries staged progression
as a sequence of attestations; the schema's five-slot structure
can be re-used by reinterpretation (boarding initiated, cabin
ready, pushback, takeoff, arrival, or whatever projection the
assembly designer adopts), without authoring a new schema.
\texttt{figaro-handoff-v1} carries the operational handoffs
between operational parties.
\texttt{figaro-jurisdiction-v1} carries the off-chain
forum-selection clause for international flights, where the
parties want to specify a particular dispute-resolution forum.
\texttt{figaro-geo-v1} carries the origin and destination
geohashes where geographic-anchoring is operationally relevant.
The GHG family of schemas is available for assemblies that wish
to compose carbon-emission attestations into the ticket
commitment.

The credential schemas verifying the wallet-held NFTs that
qualify each provider to participate (the pilot wallet's FAA
certification, the aircraft wallet's airworthiness certificate,
the maintenance-provider wallet's authority-issued certifications)
are runtime-tier credential schemas the assembly designer
composes; the kernel does not enforce them, but the validator
contracts the assembly binds at agreement signing do.

\paragraph{The new sister-schema pair.}
The pair binds a commerce-v1 ticket commitment to a specific
scheduled departure event and a specific scheduled arrival event,
so the resolution path can verify whether the sellers delivered
against schedule at each end of the journey. Splitting departure
and arrival into a sister pair follows the established
sister-schema pattern in Figaro (proximity-policy /
proximity-proof; the GHG-disclosure family / GHG-measurement
split): each event is attested separately, with content shaped
to the event the parties expect at that point.

\texttt{scheduled-departure-v1} (Layer~A spec):
\begin{itemize}
  \item \texttt{flightDesignator} (string, IATA format, e.g.\ ``NK1234'')
  \item \texttt{scheduledDeparture} (ISO~8601 datetime, with explicit timezone)
  \item \texttt{originAirport} (string, ICAO~4-character, e.g.\ ``KORD'')
  \item \texttt{aircraftType} (string, optional, IATA aircraft code; omitted
    if the parties reserve substitution rights at the runtime tier)
  \item \texttt{departureToleranceSeconds} (integer, the agreed slip below
    which no penalty applies; e.g.\ 900 for a 15-minute window)
\end{itemize}

\texttt{scheduled-arrival-v1} (Layer~A spec):
\begin{itemize}
  \item \texttt{flightDesignator} (string, IATA format)
  \item \texttt{scheduledArrival} (ISO~8601 datetime, with explicit timezone)
  \item \texttt{destinationAirport} (string, ICAO~4-character)
  \item \texttt{aircraftType} (string, optional, IATA aircraft code)
  \item \texttt{arrivalToleranceSeconds} (integer)
\end{itemize}

The two schemas share the \texttt{flightDesignator} field so the
arrival event can be matched to its departure within the bonded
process. We deliberately do not impose a maximum on
\texttt{departureToleranceSeconds} or
\texttt{arrivalToleranceSeconds}. Scheduled-service operations
vary widely in operationally meaningful tolerance windows, and
hard-coding a maximum at the validator contract would foreclose
legitimate use cases under the first-write-wins binding discipline
that makes the schemaId permanent. The contracting parties choose
the tolerance their operation supports; if a domain-specific
maximum is later wanted, it can be enforced at a tighter sister
schema (e.g.\ \texttt{figaro-passenger-departure-v1}) without
mutating the present binding. The pair is the seed of a
\emph{scheduled-service} family; later extensions could cover
scheduled ground transport, scheduled deliveries with tight
windows, and other scheduled-binding use cases without
reauthoring the family.

\paragraph{Schemas the architecture deliberately does not include.}
Three classes of schema are conspicuously absent. First, there is
no force-majeure schema that allows a counterparty to avoid bond
loss in weather events. Force-majeure handling is a
\emph{composition} concern, not a kernel concern: a party who
wants to insure against weather-driven delay composes a parallel
insurance process; the kernel neither knows about weather nor
adjudicates whether a given event qualifies. This is the
no-escape-hatches discipline of the kernel, applied to the
air-service domain. Second, there is no schema that governs which
delays count toward bond loss; the
\texttt{scheduled-departure-v1} and
\texttt{scheduled-arrival-v1} tolerances are the kernel's only
such criteria, and any further refinement is for the off-chain
forum that adjudicates disputes the parties bring to it. Third,
there is no schema that grants any party a unilateral reschedule
right; rescheduling is a new commitment with a new schemaId or a
new instance, not a mutation of the existing binding. Append-only
identity is the discipline.

% ════════════════════════════════════════════════════════════════════
\section{Comparison with the Conventional Apparatus}
\label{sec:comparison}

The architectural change should be located precisely against the
existing apparatus to clarify what does and does not change.

\paragraph{What changes.}
The architecture changes the cost-flow path. Under the
contract-of-carriage regime, passenger schedule loss is bounded
by per-incident penalty caps and counterparty discretion in
issuing voucher compensation; sub-supplier liability for
sub-service slip is bounded by per-incident contractual penalties;
and the residual loss is absorbed by whichever contract has the
loosest cap. Under the bonded architecture, schedule loss is
allocated proportionally to the wallet whose bond is exposed at
the relevant process layer: sub-supplier bond for sub-service
slip, root-counterparty bond for arrival-tolerance violation,
passenger bond returned in full on satisfactory delivery. The
cost flow is direct, proportional, and tracks the operational
layer at which performance failed.

\paragraph{What does not change.}
Safety regulation is unchanged: the FAA in the US, EASA in the
EU, and analogous national civil-aviation authorities continue to
certify aircraft, pilot licensing, maintenance procedures, and
operational standards. Crew-labor protections are unchanged:
pilot-union and flight-attendant collective bargaining and
flight-time / duty-time limits under FAR Part 117 in the US (with
analogous regimes elsewhere) apply to the crew-edge under the
bonded architecture exactly as they apply under the conventional
regime. Insurance and reinsurance are unchanged in their
underlying function; the bonded architecture changes the risk
profile aviation insurance markets price (a smaller residual
after bond allocation; concentrated risk at whichever edges carry
the largest cumulative-value bond posture; a new class of
bond-default risk at every provider wallet in the process) but
does not displace the insurance function. Antitrust and
competition law are unchanged: provider-wallet markets remain
subject to anti-collusion enforcement, slot-allocation regimes at
congested airports remain governed by the relevant authorities,
and code-share arrangements remain subject to the applicable
competition regimes.

\paragraph{Regulatory roles attached to wallets.}
Some regulatory functions attach to a single designated holder
in the conventional regime --- FAA Part 121 (and analogous regimes
in other jurisdictions) imposes a single-certificate-holder
requirement on commercial scheduled passenger operations, with
many safety functions attaching to that holder. Under the bonded
architecture the certificate-holder is one of the wallets in the
assembly (typically the aircraft wallet, or a wallet associated
with the aircraft's operating credentials); the certificate
function is preserved, but the certificate-holder participates as
one bonded counterparty alongside others, not as a cross-market
coordinator. Off-chain discovery surfaces, quality-signal
aggregators, and multi-process curators may also be useful around
the bonded process. The kernel forbids any party other than the
passenger from being the buyer of resource-market commits, so no
off-chain surface can re-emerge as a kernel-level coordinator.

\paragraph{What this is not.}
The architecture is a re-architecture of the cost-flow path in
air-service coordination. Whether any specific provider wallet
operates well or poorly within the architecture is an
operational-management question the architecture is silent on.

The architecture is also not a passenger-protection regime. It is
a coordination protocol that has the side effect of routing
schedule-loss damages to the passenger more proportionally than
the contract-of-carriage regime does. Passenger-protection
regimes (the EU's Regulation 261/2004, the US's Department of
Transportation refund and compensation rules, and analogous
regimes elsewhere) continue to operate in their own register and
apply to the regulated wallet in the relevant jurisdiction. The
bonded architecture's passenger-side benefit is structural rather
than regulatory.

% ════════════════════════════════════════════════════════════════════
\section{Scope}
\label{sec:scope}

The argument bounds itself by several scope conditions.

\paragraph{Asset specificity.}
The architecture works best where provider wallets compete in
markets with multiple available alternatives. Where a provider is
a near-monopolist (as airport authority wallets often are at hub
airports for gate operations, or as ATC services are by
regulatory construction), the buyer-side leverage at
sub-procurement is constrained. The architecture is most effective
in the low-specificity region of the resource portfolio and least
effective at the high-specificity end.

\paragraph{Bond-currency access.}
The architecture presumes that all parties can post collateral in
a denomination the kernel accepts. For some provider wallets
(e.g.\ small catering contractors operating in jurisdictions with
limited crypto on-ramps) this may be a binding operational
constraint. The constraint is real but solvable through standard
treasury and on-ramp infrastructure.

\paragraph{Aggregate liquidity at the provider tier.}
A high-volume provider wallet covering thousands of flights per
day under the bonded architecture would face aggregate
bond-custody requirements that scale with daily volume. This is
an operational treasury question: the wallet maintains a working
bond pool sized to its active commitments, with bonds released
and rebound as processes resolve. The cumulative custody at any
moment is bounded by the active-process count multiplied by
per-process bond requirements; the architecture is
liquidity-intensive in steady state but not unbounded.

\paragraph{Composability with the conventional regime.}
The architecture is composable with the conventional regime: a
provider wallet can serve some passengers under the bonded
architecture and others under the contract-of-carriage regime,
and a sub-supplier can serve some counterparties under bonded
contracts and others under conventional ones.

% ════════════════════════════════════════════════════════════════════
\section{Conclusion}
\label{sec:conclusion}

Air service is a coordination problem across resource markets.
Each resource provider is a wallet holding productive value:
NFT credentials that qualify the wallet to participate, the
underlying real-world asset whose lifecycle the wallet tracks,
and a signing key that binds the wallet to bonded commitments.
The bonded primitive in \texttt{FigaroCore} coordinates across
those wallets without a third party: each commit binds a
bilateral edge, asymmetric bonding makes cooperation the
dominant strategy at every edge, and atomic resolution propagates
a single commit-or-don't decision through the entire process.
Cascading-delay reduction follows as a consequence of bond
posture and atomic resolution; it is a property the architecture
produces, not the property the paper is for.

The architecture is composable with existing safety regulation,
crew-labor protections, insurance markets, and competition
regimes. It changes the cost-flow path without displacing the
regulatory and contractual layers that surround it. The
architecture is well-defined, and its operational properties
differ from the conventional regime in ways the operations-research
literature on airline scheduling has documented without being
positioned to address. What the present paper supplies is the
assembly --- the configuration of bonded commits, the credential
schemas, the schedule sister-pair, and the disruption-resolution
mechanism --- that turns the bonded primitive into a working
air-service coordination architecture.

\section*{Acknowledgments}
The worked example originated in a project repository of
agent-coordination scenarios; the present paper formalizes the
scenario into an industrial-engineering register.

% ════════════════════════════════════════════════════════════════════
\begin{thebibliography}{99}

\bibitem[Lan et al.(2006)]{lan_clarke_barnhart_2006}
Lan, S., Clarke, J.-P., and Barnhart, C.
\newblock Planning for Robust Airline Operations: Optimizing Aircraft
  Routings and Flight Departure Times to Minimize Passenger Disruptions.
\newblock \emph{Transportation Science}, 40(1):15--28, 2006.

\bibitem[European Parliament and Council(2004)]{eu261_2004}
European Parliament and Council.
\newblock \emph{Regulation (EC) No 261/2004 establishing common rules
  on compensation and assistance to passengers in the event of denied
  boarding and of cancellation or long delay of flights}.
\newblock February 11, 2004.

\bibitem[FAA(2014)]{far117}
US Federal Aviation Administration.
\newblock \emph{14 CFR Part 117: Flight and Duty Limitations and Rest
  Requirements: Flightcrew Members}.
\newblock 2014.

\end{thebibliography}

\end{document}
